\documentclass[a4paper,11pt]{article}
\usepackage[T1]{fontenc}
\usepackage[utf8]{inputenc}
\usepackage{lmodern}
\usepackage{microtype}
\usepackage[a4paper,margin=1in,headheight=15pt,headsep=14pt]{geometry}
\usepackage{amsmath,amssymb}
\usepackage{graphicx}
\usepackage{booktabs}
\usepackage[table,svgnames]{xcolor}
\usepackage[most]{tcolorbox}
\tcbuselibrary{skins,breakable}
\usepackage{pifont}
\IfFileExists{bbding.sty}{%
  \usepackage{bbding}%
  \newcommand{\glyphHand}{\Pointinghand}%
  \newcommand{\glyphPencil}{\Pencil}%
}{%
  \newcommand{\glyphHand}{\ding{43}}%
  \newcommand{\glyphPencil}{\ding{46}}%
}
\usepackage{enumitem}
\setlist{leftmargin=1.5em,topsep=4pt,itemsep=3pt}
\newlist{steps}{enumerate}{1}
\setlist[steps]{label=\textbf{Step \arabic*.},
  leftmargin=4.4em, labelwidth=3.8em, labelsep=0.6em, labelindent=0pt,
  align=left, topsep=5pt, itemsep=6pt}
\usepackage{parskip}
\usepackage{fancyhdr}
\usepackage{titlesec}
\usepackage[hidelinks]{hyperref}
\usepackage{xurl}
\hypersetup{breaklinks=true}

\definecolor{navy}{HTML}{21355E}
\definecolor{navyrule}{HTML}{21355E}
\definecolor{inktext}{HTML}{1A202C}
\definecolor{mutedtext}{HTML}{6B7280}

\definecolor{intuFrame}{HTML}{2C5AA0}   \definecolor{intuFill}{HTML}{F4F7FC}
\definecolor{everyFrame}{HTML}{8A5A1E}  \definecolor{everyFill}{HTML}{FBF1E3}
\definecolor{workFrame}{HTML}{2E7D52}   \definecolor{workFill}{HTML}{F1F8F3}
\definecolor{watchFrame}{HTML}{B23A48}  \definecolor{watchFill}{HTML}{FBF1F2}
\definecolor{keyFrame}{HTML}{6A4C93}    \definecolor{keyFill}{HTML}{F4F0F9}
\definecolor{waitFrame}{HTML}{0F766E}   \definecolor{waitFill}{HTML}{EEF7F5}

\color{inktext}
\hypersetup{linkcolor=navy,urlcolor=navy,citecolor=navy}

\tcbset{
  housecallout/.style 2 args={
    enhanced, breakable,
    colback=#2, colframe=#1,
    boxrule=0.6pt, arc=2.4pt, outer arc=2.4pt,
    left=10pt, right=10pt, top=9pt, bottom=9pt,
    before skip=11pt, after skip=11pt,
    fontupper=\normalsize,
    coltitle=white,
    fonttitle=\bfseries\sffamily\small,
    attach boxed title to top left={xshift=6mm,yshift=-3mm},
    boxed title style={
      colback=#1, colframe=#1,
      rounded corners, arc=3pt,
      boxrule=0pt, left=6pt, right=6pt, top=2pt, bottom=2pt
    }
  }
}

\newtcolorbox{intuition}{housecallout={intuFrame}{intuFill}, title={\raisebox{-0.05ex}{\glyphHand}\,\ The intuition}}
\newtcolorbox{everyday}{housecallout={everyFrame}{everyFill}, title={\raisebox{-0.05ex}{$\bigstar$}\,\ Everyday picture}}
\newtcolorbox{worked}[1]{housecallout={workFrame}{workFill}, title={\raisebox{-0.05ex}{\glyphPencil}\,\ Worked example: #1}}
\newtcolorbox{watchout}{housecallout={watchFrame}{watchFill}, title={\raisebox{-0.05ex}{\ding{55}}\,\ Watch out}}
\newtcolorbox{keytake}{housecallout={keyFrame}{keyFill}, title={\raisebox{-0.05ex}{\ding{51}}\,\ Key takeaway}}
\newtcolorbox{waitwhat}{housecallout={waitFrame}{waitFill}, title={\raisebox{-0.05ex}{\ding{93}}\,\ Wait --- really?}}

\newcommand{\CourseShort}{DNN Companion}
\newcommand{\SessionNo}{14}
\newcommand{\LectureTitle}{Unsupervised Learning \& GMM}

\pagestyle{fancy}
\fancyhf{}
\fancyhead[L]{\footnotesize\color{mutedtext}\textsf{\CourseShort}}
\fancyhead[R]{\footnotesize\color{mutedtext}\textsf{Session \SessionNo\ \textperiodcentered\ \LectureTitle}}
\fancyfoot[L]{\footnotesize\color{mutedtext}\textsf{WILP, BITS Pilani}}
\fancyfoot[C]{\footnotesize\color{mutedtext}\thepage}
\renewcommand{\headrulewidth}{0.4pt}
\renewcommand{\footrulewidth}{0pt}

\titleformat{\section}{\sffamily\Large\bfseries\color{navy}}{\thesection}{0.8em}{}[{\vspace{2pt}\color{navy!70}\titlerule[0.8pt]}]
\titleformat{\subsection}{\sffamily\large\bfseries\color{navy!92}}{\thesubsection}{0.7em}{}

\newcommand{\secsub}[1]{\noindent{\sffamily\bfseries\itshape\color{navy!85}\large #1}\par\medskip}
\newcommand{\flipanswer}[1]{\par\noindent\textbf{Answer:} #1}
\newcommand{\housefig}[2]{\begin{figure}[htbp]\centering\includegraphics[width=\linewidth]{#1}\caption{#2}\end{figure}}

\begin{document}

\thispagestyle{fancy}

\begin{tcolorbox}[enhanced, breakable=false, colback=navy, colframe=navy, boxrule=0pt, arc=3pt, outer arc=3pt, left=16pt, right=16pt, top=16pt, bottom=16pt, before skip=2pt, after skip=14pt, width=\linewidth]
  {\sffamily\bfseries\color{white}\fontsize{12.5}{15}\selectfont WILP, BITS Pilani}\par\vspace{9pt}
  {\sffamily\footnotesize\bfseries\color{white!85}\MakeUppercase{Deep Neural Networks \textperiodcentered\ Companion Reader}}\par\vspace{6pt}
  {\sffamily\bfseries\fontsize{26}{30}\selectfont\color{white} Session 14: Unsupervised Learning \& GMM\par}
  \vspace{5pt}
  {\sffamily\large\color{white!88} K-Means, Mixture Models, and Expectation-Maximization (EM)\par}
  \vspace{9pt}
  {\color{white!55}\rule{\linewidth}{0.4pt}}\par\vspace{6pt}
  {\sffamily\footnotesize\color{white!82} Read alongside the lecture slides \textperiodcentered\ Every slide example worked out in full\par}
\end{tcolorbox}

\section{Introduction to Unsupervised Learning}
\secsub{Discovering hidden patterns and structures in unlabeled data without target ground truth.}

In supervised learning, every training sample comes with a known target label $y$. In \textbf{unsupervised learning}, we receive only feature vectors $x_1, x_2, \dots, x_N$ without ground truth labels. The goal is to discover latent structure, underlying clusters, or generative probability distributions.

\begin{intuition}
Unsupervised learning is like sorting a large pile of mixed coins into groups based on size and weight without knowing coin names.
\end{intuition}

\section{K-Means Clustering}
\secsub{Partitioning $N$ samples into $K$ disjoint clusters by minimizing within-cluster squared distances.}

K-Means is a hard-clustering algorithm that partitions data into $K$ distinct clusters centered at means $\mu_1, \dots, \mu_K$.

\subsection{Objective Function and Algorithm}
K-Means minimizes the distortion measure $J$:
\begin{equation}
J = \sum_{n=1}^{N} \sum_{k=1}^{K} r_{nk} \|x_n - \mu_k\|^2
\end{equation}
where $r_{nk} = 1$ if sample $x_n$ is assigned to cluster $k$, and $r_{nk} = 0$ otherwise.

\begin{enumerate}[leftmargin=*]
    \item \textbf{E-Step (Assignment):} Assign sample $x_n$ to the nearest centroid $\mu_k$:
    \begin{equation}
    r_{nk} = \begin{cases} 1 & \text{if } k = \arg\min_j \|x_n - \mu_j\|^2 \\ 0 & \text{otherwise} \end{cases}
    \end{equation}
    \item \textbf{M-Step (Update):} Recompute centroids by averaging assigned points:
    \begin{equation}
    \mu_k = \frac{\sum_{n=1}^N r_{nk} x_n}{\sum_{n=1}^N r_{nk}}
    \end{equation}
\end{enumerate}

\housefig{figures/fig_kmeans_vs_gmm.png}{Comparison of hard K-Means boundaries versus soft Gaussian Mixture Model probabilistic boundaries.}

\section{Gaussian Mixture Models (GMM)}
\secsub{Soft probabilistic clustering representing data as a weighted mixture of $K$ Gaussian distributions.}

While K-Means makes hard binary assignments, a \textbf{Gaussian Mixture Model (GMM)} calculates the soft probability (responsibility $\gamma_{nk}$) that point $x_n$ belongs to Gaussian component $k$.

\subsection{GMM Formulation}
The overall probability density function is a linear combination of $K$ Gaussians:
\begin{equation}
p(x) = \sum_{k=1}^{K} \pi_k \mathcal{N}(x \mid \mu_k, \Sigma_k)
\end{equation}
where $\pi_k$ are mixture weights satisfying $\sum_{k=1}^K \pi_k = 1$ and $\pi_k \ge 0$.

\section{The Expectation-Maximization (EM) Algorithm}
\secsub{Iterative optimization framework for finding maximum likelihood parameter estimates in latent variable models.}

\subsection{EM Algorithm Steps}
\begin{steps}
\item \textbf{Initialize:} Set initial parameters $\mu_k, \Sigma_k, \pi_k$.
\item \textbf{Expectation Step (E-Step):} Calculate responsibilities $\gamma_{nk}$:
\begin{equation}
\gamma_{nk} = \frac{\pi_k \mathcal{N}(x_n \mid \mu_k, \Sigma_k)}{\sum_{j=1}^K \pi_j \mathcal{N}(x_n \mid \mu_j, \Sigma_j)}
\end{equation}
\item \textbf{Maximization Step (M-Step):} Re-estimate parameters using $\gamma_{nk}$:
\begin{equation}
N_k = \sum_{n=1}^N \gamma_{nk}, \quad \mu_k^{\text{new}} = \frac{1}{N_k} \sum_{n=1}^N \gamma_{nk} x_n, \quad \pi_k^{\text{new}} = \frac{N_k}{N}
\end{equation}
\end{steps}

\housefig{figures/fig_em_convergence.png}{The EM algorithm monotonically increases the log-likelihood $L(\theta)$ until convergence.}

\begin{worked}{GMM Numerical Calculation}
Let us execute the GMM numerical slide example step-by-step:

Consider 1D data points $X = [1, 2, 5, 6]$. We fit $K=2$ Gaussians.

Initial parameters: $\mu_1 = 1.5$, $\mu_2 = 5.5$, variances $\sigma_1^2 = 1$, $\sigma_2^2 = 1$, weights $\pi_1 = 0.5, \pi_2 = 0.5$.

\begin{steps}
\item \textbf{E-Step for $x_1 = 1$:}
Compute $\mathcal{N}(1 \mid 1.5, 1) = \frac{1}{\sqrt{2\pi}} e^{-0.5(1-1.5)^2} = 0.352$.
Compute $\mathcal{N}(1 \mid 5.5, 1) = \frac{1}{\sqrt{2\pi}} e^{-0.5(1-5.5)^2} = 0.000014$.
Responsibility $\gamma_{1,1} = \frac{0.5 \times 0.352}{0.5 \times 0.352 + 0.5 \times 0.000014} \approx 0.9999$.
\item \textbf{M-Step Update for $\mu_1$:}
Update mean $\mu_1^{\text{new}} = \frac{\sum_n \gamma_{n1} x_n}{\sum_n \gamma_{n1}} \approx \frac{1(1) + 1(2) + 0(5) + 0(6)}{1 + 1 + 0 + 0} = 1.5$.
\end{steps}
\end{worked}

\section{Summary \& Self-Test}
\begin{keytake}
K-Means performs hard clustering by minimizing squared distance. GMM extends this to soft probabilistic clustering using the EM algorithm, updating responsibilities (E-step) and Gaussian parameters (M-step).
\end{keytake}

\subsection{Self-Test Quiz}
\textbf{Q1:} What is the main difference between K-Means and GMM?

\flipanswer{K-Means makes hard assignments, whereas GMM computes soft probabilistic responsibilities.}

\end{document}
